\documentclass[12pt]{article}

\usepackage{amssymb}
\usepackage{amsthm}
\usepackage{amsfonts}
\usepackage{booktabs}
\usepackage{commath} 
\usepackage{dcolumn}
\usepackage{graphicx}
\usepackage{hyperref}
\usepackage{natbib}
\usepackage{setspace} 
\usepackage{subcaption}
\usepackage{dsfont}

\newcommand{\Lsm}{$\mathcal{L}$}
\newcommand{\Hsm}{$\mathcal{H}$}

\begin{document}

\section{Model}

\subsection{Motivation}
The experimental intervention only makes sense in market with substantial search frictions---if the two sides could easily find the counter-party of their choice, no additional effort to help direct search would be needed.

In the directed search literature, by making \emph{ex ante} commitments and announcements, buyers/workers can direct their search.
There are still matching frictions, but they happen ``later'' after parties have used ex ante information (like announced prices) to decide where to search.
In the directed search literature, there is often the notion of sub-markets (say all buyers announcing a given price that they are willing to pay), with sellers then joining sub-markets in response to these announcements, weighing off the price they will pay with their matching probability.
In these models, these sub-markets are often conjured only to show they will not actuall exist---in the sense that the big so-what is that that separate sub-markets cannot exist in equilibrium (and hence price dispersion cannot exist in equilibrium). 

Our situation is different from the rest of the literature, in that sub-markets existing in equilibrium.  
The reason we might try to induce these sub-markets is to improve something about matching. 
By creating sub-markets, we potentially thin out the markets and reduce the total quantity of ``meetings,'' but at the benefit that both sides are more likely to find a suitable counter-party.
Assuming there are gains from creating sub-markets, there is next the question of whether that separation can actually be implemented.
In particular, we need to check whether all workers and firms will find it incentive compatible to participate in ``their'' sub-market. 

\subsection{Worker and firm types}
Consider a labor market with types of workers, H and L, with the fraction of H-type workers being $g$ (think ``g'' for ``good'').
There are also two types of firms, H and L, with the fraction of H-type firm being $p$ (think ``p'' for ``picky'').
The market is balanced in the sense that there is a unit measure of both firms and workers. 

\subsection{Production} 
Firms need to hire a single worker to produce the output.
To produce the output, firms need workers that is their type or higher:
an L-type firm can produce the output with an L-type worker or an H-type worker, but an H-type firm needs an H-type worker.
The value of the output for an L-type firm is $v_L$ and for an H-type firm, $v_H$, with $v_H > v_L$.
It is costless for the worker to produce the output, but the firm does have to pay wages. 
For now, we will set aside how wages are determined---and hence pay-offs are determined---and simply state wages exist.

\subsection{Meetings and hiring}
Neither the firm nor the worker can directly find the counter-party of the compatible type.
They first have to ``meet'' according to a balanced constant returns to scale matching function, $m(n_b, n_s) = \mu \sqrt{n_b n_s}$, that operates at the level of the market.
If a firm and a worker meet, the firm observes the worker's type and decides whether to hire the worker at their proposed wage.
We use the term ``meet'' rather than ``match'' to emphasize that the matching function only returns the \emph{opportunity} to hire, rather than a hire.
Hiring is a two-step process---first the firm and worker need to meet, then the firm needs to decide they want to hire that worker they met.

With all workers and firms participating in the market---which we call the ``pooled'' market---the number of meetings that occur is $m(1,1) = \mu$, but the number of hires is only
\begin{align*}
\mbox{Hires} = \left( pg + (1 - p) \right)\mu.
\end{align*}
Among the meetings, the number of hires is comprised of (a) the picky firms that happen to meet with good workers ($pg$) and (b) the non-picky firms, who hire anyone they meet with ($1-p$).
To emphasize, not all meetings lead to a hire---when H-type firms meets with L-type workers, they decide not hire.
The hires that do no occur are the $(1-g)p$.
The total surplus from the pooled market is 
\begin{align}
  S = \left(v_H pg + v_L (1 - p) \right)\mu.
\end{align}

Aside from the hires that do not occurr, another defect of this pooled equilibrium is some hires are of poor ``quality.''
In particular, there are H-type workers hired by L-type firms.
This is better than no hire at all, but it is not as good as if those H-type workers could have met with H-type firms. 

%Let $w_L$ be the wage paid by L-type employers and $w_H$ be the wage paid by H-type employers.
%If there is a single market-clearing wage, then $w_L = w_H$.
%If there is bargaining, then we could have, even in the decentralized equilibrium, different wages (e.g., $w_L = \frac{1}{2} v_L$ and $w_H = \frac{1}{2} w_H$). 

%% \subsection{Wages} 
%% It is costless for both sides to participate in the matching market, but firms have to pay a cost $k$ to hire.
%% Firms tell an auctioneer the most they would be willing to pay to hire a worker that can produce the output.
%% The market wage is the lowest bid submitted and prevails for all relationships.
%% This is $v_L - k$.

\subsection{Enabling directed search by creating sub-markets}
Now consider a market structure where the auctioneer can separate firms and workers into two sub-markets, \Lsm{} and \Hsm{}.
Meetings also occurr within these sub-markets according to the same matching function. 
The equilibrium we are interested in is a pure equilibrium where all market participants go to the market of their type:
we have an \Lsm{} sub-market where all the firms and workers are L-types and an \Hsm{} sub-market where all the firms and workers are H-types.

\verb+JJH: What would happen if we allowed mixed strategies?+

The benefit of this separation is that all meetings lead to a hire and the hires are of the best possible quality:
sorting firms and workers never meet with a counter-party that is not their type.
The downside of this sorting is that there are fewer meetings---and hence chances to hire---except in the edge case of perfect balance (i.e., $p = g$). 
The reason for the loss in meetings is that:
\begin{align*}
  m(1,1) \ge m(p,g) + m(1-p, 1-g). 
\end{align*}

The market designer trade-off is fairly straightforward:
by creating sub-markets, we almost certainly reduce the number of meetings, but we might increase the number and quality of hires.
If we expect firms and workers to sort into the right sub-market, we have to consider incentive compatibility.
Will firms and workers try to go the right market for their type?

\subsection{Worker incentive compatibility}
On the worker side, first note that we do no need to worry about incentive compatibility for the L-type worker.
They would never want to choose the \Hsm{} sub-market---if they happen to get a meeting, they will be revealed as an L-type to the H-type firm, which will not hire them.
The case to consider is the H-type worker. 

In the \Hsm{} sub-market, there will be $m(p,g)$ meetings, and so an H-type worker's probability of being hired is $m(p,g)/g$.
In the \Lsm{} sub-market, there will be $m(1 - p, 1 - g)$ meetings, and so the H-type worker's probability of being hired if they switch to the \Lsm{} sub-market is $m(1-p,1-g)/(1-g)$.
As such, for the H-type worker, incentive compatibility requires that 
\begin{align*}
  \frac{m(p,g)}{g} w_\mathcal{H} \ge \frac{m(1 - p, 1 - g)}{1 - g}w_\mathcal{L}.
\end{align*}

From this formula, we can gain some intuition for the comparative statics of the equilibrium. 
For a given $p$, as the number of other H-type workers increases (higher $g$), the \Hsm{} market becomes less attractive.
The problem from a H-type worker's perspective is that although more H-type workers increase the total number of meetings, those meetings are split over a larger pool of workers, and so meeting/hiring probability goes down, as $\frac{\partial}{\partial g}[\frac{m(p,g)}{g}] < 0$.
%For the same reason, a higher $g$ makes the L market relatively more attractive to the H worker. 
For a given $g$, more picky firms (higher $p$) makes the \Hsm{} market more attractive to an H-type worker---there are more meetings but they still split over the same fixed number of H-type workers.
Unsurprisingly, all else equal, the higher the wage in \Hsm{} sub-market and the lower the wage in \Lsm{} sub-market, the more attractive the H sub-market is to H-type workers. 

The incentive compatibility constraint for the H-type worker binds for a $\hat{p}_S$ (think ``S'' for supply-side) when
\begin{align*}
  \hat{p}_S = \frac{1}{\left(\frac{1}{g} - 1\right) \left(\frac{w_H}{w_L}\right)^2 + 1}.
\end{align*}
With $g$ is higher (more H-workers), more H-type firms are needed to keep the H-type worker indifferent: $\frac{\partial \hat{p}}{\partial g} > 0$. 
If wages in the H sub-market are higher than the L sub-market, then a relatively worse meeting probability can keep the H-type worker indifferent: $\frac{\partial \hat{p}}{\partial (w_H/w_L)} < 0$.  

\subsection{Firm incentive compatibility}
Mirroring the worker case, on the firm side, we only have to consider one firm type---the L-type firm. 
L-type firms can produce the output with either kind of worker, whereas the H-type firms are not interested in the \Lsm{} sub-market, as they do not want to hire any worker they would meet with.

Let $\pi_{\mathcal{L}|L} = v_L - w_{\mathcal{L}}$ be a L-type firm's profits from hiring in the \Lsm{} sub-market and paying $w_\mathcal{L}$ in wages. 
If they instead join the \Hsm{} sub-market, their profit would be $\pi_{\mathcal{H}|L} = v_L - w_{\mathcal{H}}$. 

In the \Hsm{} sub-market, there will be $m(p, g)$ meetings, and so a firm's probability of hiring is $m(p, g)/p$.
In the \Lsm{} sub-market, there will be $m(1 - p, 1 - g)$ meetings, and so the firm's probability of hiring is $m(1 - p, 1 - g)/(1 - p)$.
For the L-type firm, incentive compatibility requires that 
\begin{align*}
  \frac{m(1 - p, 1 - g)}{1 - p} \pi_{\mathcal{L}|L} \ge \frac{m(p, g)}{p} \pi_{\mathcal{H}|L}
\end{align*}
The incentive compatibility constraint binds for a $\hat{p}_D$ (think ``D'' for demand-side), given $g$ when
\begin{align*}
 \hat{p}_D = \frac{1}{\left(\frac{1}{g} - 1\right) \left(\frac{ \pi_{\mathcal{L}|L}}{\pi_{\mathcal{H}|L}} \right)^2 + 1}.
\end{align*}

For the L-type firm, they only get $v_L$ regardless of which worker they hire, but they have to pay higher wages in the \Hsm{} sub-market.
As such, $\pi_{\mathcal{L}|L} > \pi_{\mathcal{H}|L}$.
Thus for the L-type firm to entertain switching, it is has to be because the \Hsm{} offers a higher probability of meeting a worker. 
If there is no difference in wages across the two sub-markets $(\pi_H = \pi_L)$, in which case $\hat{p} = g$, and for a separating equilibrium to be possible, $p > g$.
Without any cost savings, the L-type firm has to find the \Hsm{} sub-market as unattractive because it offers a low meeting probability---which in turn is the result of lots of other firms in \Hsm{} and relatively few workers.  

\subsection{Which constraint binds?}

There are two incentive compatibility constraints with similar structures.
For a given $g$, there are two threshold picky firm fractions---$\hat{p}_D$ and $\hat{p}_S$ needed to support a separating equilibrium.
We are interested in $\max\{\hat{p}_D, \hat{p}_S\}$.

Two ratios matter for determining which incentive compatibility constraint matters---$\pi_L / \pi_H$ and $w_H / w_L$.
The smaller ratio is the one that determines the binding constraint. 
If $\pi_L / \pi_H > w_H / w_L$, then $w_H / w_L$ determines the $\hat{p}$ and it is L-type worker we have to consider.

Let the ratio of wages be $r = w_H / w_L$ and let the ratio of L-type firm value to L sub-market wage be $s_L = v_L/w_L$.
The profit ratio is greater than the wage ratio when 
\begin{align*}
  \frac{\pi_L}{\pi_H} & = \frac{v_L - w_L}{v_L - w_H} > \frac{w_H}{w_L} \\
                      &  = \frac{v_L/w_L - 1}{v_L/w_L - w_H/w_L} > \frac{w_H}{w_L} \\
                      & =  \frac{s_L - 1}{s_L - r} > r.
\end{align*}
If $s_L < 2$, this holds when  $s_L > r$.
If $s > 2$, then it holds if $s_L > r$ and $s_L < r + 1$.

If $w_L$ is a really small share of $v_L$, making $s_L$ it is going to be an even smaller share of $v_H$.

What if Nash bargaining?
In this case, $w_L = \frac{1}{2} v_L$ and  $w_H = \frac{1}{2} v_H$.
\begin{align*} 
  \frac{\pi_L}{\pi_H} & = \frac{v_L - w_L}{v_L - w_H} > \frac{w_H}{w_L} \\
  & = \frac{v_L}{2 v_L - v_H}     > \frac{v_H}{v_L} \\
  & = \frac{1}{2  - r}   >  r
\end{align*}
Which is true if $r > 1$. 

\subsection{Social welfare from creating sub-markets}
The social surplus with two sub-markets is
\begin{align*}
 S = v_H m(p,g) + v_L m(1-p, 1 - g)
\end{align*}
and so the sub-market approach is more efficient than the pooled equilibrium when:
\begin{align*}
  \Delta S = v_H \left(m(p,g) - pg \mu  \right)  + v_L \left( m(1 - p, 1 - g) - (1 - p) \mu \right) > 0.
\end{align*}

Note that as both $p < 1$ and $g < 1$, $\frac{1}{\mu} m(p,g) > pg$, and so we always have more high-type \emph{hires} from creating sub-markets. 
As we might intuit, all else equal, a greater social value from creating H-type hires recommends market-splitting. 

Although H-type hires always increase, the effects on the number of low-type hires is ambiguous. 
There are more low-type matches with a sub-market when $\frac{1}{\mu} m(1 - p, 1 - g) - (1 - p) > 0$.
This is the case when $g < p$, i.e., when there are relatively more picky firms than H-type workers.
In this case, market-splitting is kind of a free lunch---we get more hires of both kinds. 
The intuition is that when there are lots of picky firms and not many H type workers, many of the ``matches'' in the pooled market led to non-hires: picky firms paired with not-good workers do not hire, and so the compensatory H-worker/L-firm matches do not occur, as is the case when the good workers are the short side of the pooled market.  
By creating sub-markets when there are lots of picky firms, otherwise matched but un-hired L-type workers are hired. 

If $g > p$, the welfare case for splitting into sub-markets is less straight-forward---the increase in $v_H$ has to weighed against the loss from fewer $v_L$ matches.
However, it is ``hard'' for $g > p$ but still have incentive compatibility unless there are large wage/profit differences---the problem is when $g$ exceeds $p$---lots of good workers and not many picky firms---the \Hsm{} has a low hire probability for H-type workers and so unless wages in \Lsm{} are really low, or wages in \Hsm{} are really high, the marginal H-type worker is inclined to defect.
Similarly, when $g > p$, for the L-type firm, jumping to the \Hsm{} sub-market is tempting because it offers a high hire probability so unless \Hsm{} wages are really high or \Lsm{} wages are low, incentive compatibility is hard.  

\subsection{Wage bids}

\paragraph{Assumig TIOLI offers} 
Consistent with our empirical context, we assume that workers have to propose wages \emph{ex ante} and that these offers are TIOLI.
In the pooled market, an L-type worker knows that if they meet an H-type firm, they will not be hired, and so they only need to consider the effect of their wage offer when they meet an L-type firm.
Given that hires can only happen with a meeting, the L-type worker knows they can extract the full surplus and so they bid $w_L = v_L$.

For an H-type worker, the situation is more complex, as they can potentially be hired by both kinds of firms. 
As they can meet with a L-type firm, if they want to be hired, they have to bid no more than $v_L$.
But this is a very low bid when paired with an H-type firm---the worker could bid upto $v_H$ and still be hired.
In the pooled market, a worker meets a H-type firm with probability $p$. 
In the pooled equilibrium, an risk-neutral worker will bid $v_H$ if $p v_H > v_L$ and $v_L$ otherwise. 
If workers differ in their risk-aversion, we might see some fraction choosing $w = v_H$ and some fraction choosing $v_L$.

If $p < v_L / v_H$, then there is a single market-clearing wage, $v_L$.

If $p > v_L / v_H$, then the average wage is $pg v_H + (1-p)(1- g)v_L$.
Note that we get fewer hires, as now H-type worker / L-type employer pairings do not go through. 

\paragraph{Conditional TIOLI offers} 
If workers can condition offers at time of meeting, then we get $w = v_L$ when matched with L-type firms and $w = v_H$ when matched with H-type firms. 

\subsection{Full information Nash bargaining}
If workers Nash bargain, they split the surplus, and so it looks like the conditional TIOLI offers. 

\subsection{Incomplete information Nash bargaining}
If workers Nash bargain but do not know the firm's type, we should probably see H-type workers agreeing to some wage $w$ that both H-type and L-type firms would accept but with the firm getting a relatively better deal with L-type firms e.g, $\frac{1}{2} v_L < w < \frac{1}{2}v_H$.
I.e., the H-type workers shades up 

\verb+Myerson and others have stuff on Nash bargaining with incomplete info - not sure what it predicts.+

\subsection{What about with sub-markets?}
Even if TIOLI offers, the H-type workers can now bid $w_H$, as they know they will get paired with H-type firms.
There is also no uncertainty about types, so the Nash bargaining also looks different.

%% \paragraph{Wages}
%% One possibility that workers must pre-commit to an hourly wage.
%% If it was a TIOLI offer, in the decentralized, pooled equilibrium, an L-type worker would just bid $v_L$.
%% An L-type firm would just accept it, and an H-type firm would not hire them at any offer anyway.
%% The H-type worker has to make a choice:
%% they can bid $w_L$, ensuring they get hired by either kind of firm they get matched with, or they can bid $v_H$, eschewing L-type matches.

%% Another way to model wages if to assume their is Nash bargaining, in which case L-firms pay $w_L = \frac{1}{2} v_L$ and H-type firms pay $w_H = \frac{1}{2} v_H$.
%% A question, however, is whether workers know what kind of firm they are paired with---an H-type worker doesn't know if they will produce $v_L$ or $v_H$ for the employer. 

\subsection{Equilibrium preference?}

\paragraph{Firms}
A H type firm in the single market hired with probability $g$.
They paid $w_L$.
With sub-markets, they not pay $w_H$ and they hire with probability $m(p,g)/p$.
The sub-market hire rate is greater than the pooled hire rate when
\begin{align}
  \frac{m(p,g)}{p} > g,
\end{align}
which holds when $\frac{1}{p} > g$, but since $p < 1$ and $g < 1$, this is always the case.
So the high-type matching probability always increases. 
The question for H type firms then becomes whether the higher matching probability is ``worth'' the lower profits. 

For the L type firm, they were willing to hire anyone they matched with, and so the hire probability was $1$.
Now, it is $\frac{m(1 - p, 1 - g)}{1 - p}$ which, with a CRS matching function, is always lower except when the market is perfectly balanced $p = g$.
They still pay $w_L$, so they are unambiguously made worse-off.

\paragraph{Workers}

L type workers have a higher probability of being hired in the sub-market, as they will not get paired with H firms that will not hire them.
\begin{align}
  \frac{m(1 - p, 1 - g)}{1 - g} > \mu (1 - p)
\end{align}

H type workers were alwars hired.
d

% Mathematica: Reduce[m[1 - p, 1 - g]/(1 - g) > (1 - p) && p > 0 && p < 1 && g > 0 && g < 1]

\subsection{Social planner solution}

\begin{enumerate}
\item Put all the picky firms in the \Hsm{}
\item Add good workers to the \Hsm{} market until total surplus is maximized 
\end{enumerate} 

\begin{align}
  \max_{0 < g* < g} v_H m(p, g^*) + v_L m(1 - p, 1 -  g^*) 
\end{align}

The optimal
\begin{align}
  g^* = \frac{1}{\left(\frac{v_L}{v_H}\right)^2 \left(\frac{1}{p} - 1\right) + 1}
\end{align}

If a social planner could direct workers, it would pair H-type firms with H-type workers until it ran out.
If it ran out of H-type workers first ($p > g$), then the social pay-off would be
\begin{align}
  v_H g + v_L (1 - p)
\end{align}

If it ran out of H-type firms first ($g > p$), then the social pay-off would be
\begin{align}
  v_H p + v_L (1 - g)
\end{align}
or
\begin{align}
 S = v_H \min\{p, g\} + v_L \max\{1-p,1-g\}
\end{align}



\end{document}

